\section{Introduction}
\label{sec:intro}
Recent work \cite{hu2024visual, zheng2025deepeyes, fan2025grit, Realunify, sarch2025grounded} has demonstrated that the incorporation of auxiliary images in the intermediate steps of chain-of-thought (CoT) can improve the visual reasoning of multimodal large language models (MLLMs)~\cite{comanici2025gemini,bai2025qwen2,wang2025internvl3,mavors, cheng2025visual, jian2025look}. 
%Empirically, such auxiliary images improve visual reasoning through the following aspects: (1) re-emphasizing the visual features of crucial image regions by cropping or drawing bounding boxes, alleviating the gradual loss of attention to image tokens during reasoning \citep{cheng2025visual, jian2025look}. (2) serving as visualizations of the intermediate steps of the visual-related thoughts, such as auxiliary lines in a geometry problem and intermediate states in a maze problem. 
A recent trend is to acquire auxiliary images by training MLLMs to predict key region coordinates \citep{zheng2025deepeyes, fu2025refocus, su2025pixel, fan2025grit}, invoking visual tools such as grounding or depth estimation models \citep{hu2024visual, sarch2025grounded, su2025openthinkimg}, or generating executable codes \citep{zhao2025pyvision, zhang2025thyme} to modify the input image. While promising, these methods are constrained by a limited set of external tools, lacking the flexible, human-like visual reasoning that arises within an internal perceptual space. 

To emulate the flexible visual reasoning of humans, we investigate training MLLMs to reason directly in the continuous latent space. Specifically, we enable MLLMs to generate latent embeddings that serve as intermediate visual thoughts beyond textual descriptions and image embeddings, thus eliminating the need for explicit auxiliary images and overcoming the rigidity of external tools. Recent studies have begun exploring latent visual reasoning strategies \citep{yang2025machine, li2025latent, pham2025multimodal}. Most existing methods simply align the generated embeddings with those of auxiliary images while applying a next-token prediction (NTP) loss on text tokens during supervised fine-tuning (SFT), and directly apply GRPO \citep{Shao2024DeepSeekMathPT} for reinforcement learning (RL), which leads to two key limitations: \textbf{(1) poor scalability}, as alignment over thousands of image tokens incurs high computational and memory costs, and using mean pooling to compress the image tokens \citep{yang2025machine} will distort detailed visual features; \textbf{(2) insufficient optimization of latent embeddings}, since the NTP objective in SFT can be easily overfit and the GRPO loss can only be computed for text tokens, the optimization of latent embeddings are ignored; Consequently, their improvements remain limited and task-specific.

Motivated by these observations, we propose \textbf{Monet}\footnote{The name \textbf{Monet} is derived from the great Impressionist artist, as our method enables MLLMs to think abstractly in the latent visual space.}, a novel training framework for multimodal latent reasoning  that trains a text-output MLLM (Qwen2.5-VL-7B \citep{bai2025qwen2}) to perform latent reasoning through SFT and RL. The SFT stage aims at equipping the model with the fundamental ability to generate and reason with latent embeddings. To tackle limitation \textbf{(1)}, instead of directly aligning latent embeddings with those of auxiliary images, we introduce dual supervision signals. First, since latent embeddings are intended to facilitate reasoning, we align the hidden representations of text tokens corresponding to crucial visual features observed from the auxiliary images, when conditioned on either auxiliary images or generated latent embeddings. Second, to preserve visual information, we use a controlled attention mask so that latent embeddings can directly attend to auxiliary image embeddings. To overcome limitation \textbf{(2)}, we optimize the alignment loss solely through latent embeddings by stopping gradients on non-latent representations. Furthermore, we propose \textbf{V}iusal-\textbf{l}atent \textbf{P}olicy \textbf{O}ptimization (\textbf{VLPO}), which computes policy gradient directly for latent embeddings by estimating their output probability.

Our contributions are summarized as follows:
\begin{enumerate}
    \item We propose \textbf{Monet-SFT}, a three-stage supervised fine-tuning framework that trains MLLMs to generate and reason with latent embeddings.
    \item We propose \textbf{VLPO}, a novel RL algorithm tailored for latent reasoning. Unlike GRPO, which targets text reasoning, VLPO incorporates latent embeddings into the total loss by computing an approximate probability for the latent embeddings collected during rollout.
    \item We identify the limitations of existing image-text interleaved datasets: the unnecessary usage and inaccuracy of auxiliary images. To address these limitations, we further propose a multi-stage data curation pipeline to construct \textbf{Monet-SFT-152K}, a high-quality dataset with image-text interleaved chain-of-thoughts (CoT) for SFT.
    \item Extensive experiments on real-world perception and reasoning benchmarks show that Monet-SFT and VLPO outperform conventional SFT + GRPO, cropping-based methods, and prior latent visual reasoning approaches. VLPO further enhances out-of-distribution (OOD) generalization on unseen tasks.
\end{enumerate}

